% Faithful conversion of ages-iii-sprint-4; source ODF IDs and text order retained.
\hypertarget{__RefHeading___Toc4823_3118176633}{}
\refstepcounter{subsection}\subsection*{4.3.5 Sprint 4}
\label{unit:ages-iii-sprint-4}
\addcontentsline{toc}{subsection}{4.3.5 Sprint 4}
% text-p-1278
\indent \originaltab{}Para a última Sprint do projeto, sugeri aos AGES IV uma mudança na forma de organizar o trabalho do time: ao invés de trabalharmos com um conjunto fixo de tasks a serem concluídas até o code freeze, passamos a estruturar a Sprint em microentregas semanais de funcionalidades. A proposta tinha como objetivo dar mais clareza às prioridades nesse momento final do projeto, em que o volume de integrações, ajustes e validações aumenta consideravelmente. Além disso, definimos um soft code freeze uma semana antes do encerramento oficial da Sprint e combinamos a realização de uma demonstração completa da aplicação na semana final, o que nos daria tempo hábil para corrigir falhas e estabilizar o sistema antes da entrega definitiva. Essa estratégia exigiu bastante organização do time, mas ajudou a dar mais previsibilidade à reta final do desenvolvimento.\par

% text-p-1279
\indent Minha atuação na Sprint 4 foi concentrada principalmente em tasks de frontend e em endpoints de integração com o backend, especialmente relacionados à tela de sistema e ao dashboard da aplicação. No backend, desenvolvi o endpoint de métricas consolidadas do dashboard e também o endpoint consolidado da tela de sistema, ambos pensados para centralizar as informações necessárias às interfaces e facilitar a integração com os componentes visuais da aplicação. Já no frontend, atuei na integração dos componentes de entrada e saída do sistema, além do desenvolvimento do componente de antena e do modal de ajuste da antena, funcionalidades diretamente ligadas à operação do leitor RFID e à interação do usuário com o ambiente monitorado pelo sistema. De forma geral, foi uma Sprint bastante voltada à consolidação das integrações entre frontend e backend, fechando lacunas importantes da solução e preparando o sistema para a apresentação final.\par

% text-p-1280
\indent Desde a Sprint 2, venho desenvolvendo diversas tasks de frontend em parceria com o Gabriel Borges, AGES I, e essa dinâmica se manteve também ao longo da Sprint 4. Como ele demonstrou bastante interesse por desenvolvimento de componentes e telas, passei a acompanhá-lo mais de perto, compartilhando algumas práticas de frontend que eu já conhecia e conduzindo com ele a leitura das issues, o planejamento da implementação e a revisão do código antes da abertura dos PRs. Foi uma experiência interessante porque, ao mesmo tempo em que eu o orientava em alguns pontos, também estava aprendendo junto com ele em uma área com a qual eu não tinha tanta familiaridade. Essa troca acabou tornando o processo de desenvolvimento mais colaborativo e leve, especialmente em uma sprint final naturalmente mais intensa.\par

% text-p-1281
\indent Além do desenvolvimento das funcionalidades, também precisei atuar em pontos de sustentação do projeto. No meio da Sprint, corrigi uma falha na pipeline relacionada à versão do SonarCloud, que estava depreciada e impedia a análise sintática do código ao não conseguir baixar corretamente as versões de Java necessárias. Embora a correção em si tenha sido simples, ela foi importante por evitar um possível gargalo em um momento em que o time já não tinha margem para perder tempo com impedimentos técnicos. Também precisei dedicar parte do meu tempo à organização de issues e dependências no backlog, já que diversas tasks estavam mal dimensionadas ou com vínculos incorretos, inclusive havendo casos de endpoints já implementados há semanas, mas ainda listados como pendentes na Sprint 4. Esse cenário acabou gerando confusão no fluxo de trabalho e atrasando parte do desenvolvimento, exigindo uma intervenção para reorganizar melhor o que realmente ainda precisava ser entregue na reta final do projeto.\par

\begin{reflection}
% text-p-1283
\indent \originaltab{}A Sprint 4 foi, para mim, uma sprint muito marcada pela ideia de fechamento e de maturidade do projeto. A decisão de trabalhar com microentregas semanais fez bastante sentido nesse contexto, porque a reta final exige um tipo de organização diferente das sprints anteriores: já não basta apenas desenvolver novas funcionalidades, é preciso priorizar bem, integrar o que foi construído, validar o comportamento do sistema como um todo e reservar espaço para corrigir falhas antes da entrega final. Ver essa estratégia funcionando me mostrou como a gestão do fluxo de trabalho influencia diretamente a qualidade do produto entregue, especialmente quando o projeto entra em sua fase mais delicada.\par

% text-p-1284
\indent Outro ponto que me marcou bastante foi a experiência de desenvolver frontend junto com o Gabriel. Mais do que apenas dividir tasks, foi um processo em que eu pude orientar um colega mais novo em algumas práticas que já conhecia e, ao mesmo tempo, aprender junto com ele em uma área que não é a minha zona de conforto. Além disso, os problemas com a pipeline e com a organização das issues reforçaram uma percepção que foi se consolidando ao longo do projeto: não basta automatizar testes, validações e deploys se o restante do processo não estiver minimamente organizado. A última sprint deixou isso muito claro para mim:\originalspace{}\originalspace{}qualidade de software depende tanto do código quanto da forma como o time estrutura, acompanha e sustenta o próprio desenvolvimento.\par

\end{reflection}
